\begin{longtable}{|p{3cm}|p{4cm}|p{6cm}|}
\caption{Daftar Perkakas Implementasi Sistem}
\label{tbl:tools} \\
\hline
\textbf{Kategori} & \textbf{Nama Tools} & \textbf{Deskripsi} \\ \hline
\endfirsthead

\multicolumn{3}{l}{Tabel \ref{tbl:tools} Daftar Perkakas Implementasi Sistem (lanjutan)} \\
\hline
\textbf{Kategori} & \textbf{Nama Tools} & \textbf{Deskripsi} \\ \hline
\endhead

\hline
\multicolumn{3}{r}{\textit{bersambung ke halaman berikutnya}} \\
\endfoot

\hline
\endlastfoot

% ======================= A. Programming Environment ==========================
\textbf{Bahasa Pemrograman \& Lingkungan Pengembangan} &
Python (v3.10+) &
Bahasa pemrograman utama yang digunakan untuk seluruh logika sistem, termasuk 
pipeline ekstraksi data, konstruksi graf, proses \textit{retrieval}, dan 
pengembangan antarmuka pengguna. Dipilih karena ekosistemnya yang matang dalam 
pemrosesan data dan pengembangan aplikasi AI. \\ \cline{2-3}

& Google Colab &
Platform notebook Python berbasis \textit{cloud} untuk eksperimen, eksekusi skrip,
dan pengujian model tanpa konfigurasi lingkungan lokal. \\ \hline

% ======================= B. Knowledge Base ==========================
\textbf{Basis Data \& Penyimpanan Pengetahuan} &
Neo4j (Community Edition) &
DBMS berbasis graf yang digunakan untuk menyimpan struktur hubungan objek 
Kubernetes dalam bentuk \textit{node} dan \textit{edge}. Mendukung kueri Cypher untuk penelusuran 
relasi dan validasi keterhubungan antar-objek. \\ \cline{2-3}

& Neo4j Python Driver &
Pustaka resmi yang memungkinkan aplikasi Python berkomunikasi langsung dengan 
Neo4j untuk proses konstruksi graf, eksekusi Cypher, dan \textit{graph traversal}. \\ \hline

% ======================= C. Artificial Intelligence Components ================
\textbf{Komponen Kecerdasan Buatan} &
OpenAI API (GPT-4o-mini) &
Digunakan untuk proses klasifikasi intensi, penafsiran pertanyaan menjadi kueri 
struktural, serta generasi kode YAML berdasarkan konteks. Berfungsi sebagai 
komponen inti dalam modul \textit{reasoning} dan penyusunan jawaban. \\ \cline{2-3}

& Sentence-Transformers (all-MiniLM-L6-v2) &
Model \textit{embedding} ringan yang digunakan untuk mengubah teks deskripsi API dan 
pertanyaan pengguna menjadi vektor numerik untuk pencarian semantik. Ukurannya 
efisien dengan performa yang tinggi. \\ \hline

% ======================= D. Validation & Utility ==============================
\textbf{Validasi \& Utilitas Data} &
PyYAML &
Pustaka parser YAML yang digunakan untuk memvalidasi keluaran LLM, khususnya 
pada tahap pengecekan sintaksis (NF-01). \\ \hline

% ======================= E. User Interface ====================================
\textbf{Antarmuka Pengguna} &
\textit{Command-Line Interface} (CLI) &
Antarmuka berbasis perintah untuk menjalankan fungsi utama sistem 
(\texttt{generate}, \texttt{explain}, \texttt{validate}, \texttt{--help}). \\ \hline

\end{longtable}
